
\[
\begin{tikzcd}[column sep=large,row sep=large]
Y' \arrow[r,"j"] \arrow[d,"g"'] & X' \arrow[d,"f"]\\
Y \arrow[r,"i"] & X.
\end{tikzcd}
\]
\uline{On
suppose que} \(X\) \uline{est quasi-compact et a un faisceau ample}. \uline{Alors}, \uline{on a la suite exacte}:
\[
0\longrightarrow K^{\bullet}(Y)\xrightarrow{u}K^{\bullet}(Y')\times K^{\bullet}(X)
\xrightarrow{v}K^{\bullet}(X'),
\]
\uline{où on pose}, \uline{pour} \(y\in K^{\bullet}(Y)\):
\[
u(y)=\bigl(g^*(y).\lambda_{-1}(F),-i_*(y)\bigr),
\]
\uline{et pour} \(y'\in K^{\bullet}(Y')\), \(x\in K^{\bullet}(X)\):
\[
v(y',x)=j_*(y')+f^*(x),
\]
\uline{les notations étant celles de} 3.1 \uline{et} 3.3. \uline{En outre}, \(u\) \uline{admet un inverse à gauche} \(u'\), \uline{défini par} \(u'(y',x)=g_*(y')\).

\uline{Si de plus} \(X\) \uline{est noethérien}, \uline{et} \(Y\) \uline{régulier}, \uline{alors} \(v\) \uline{est surjectif}, \uline{donc on a la suite exacte}:
\[
0\longrightarrow K^{\bullet}(Y)\xrightarrow{u}K^{\bullet}(Y')\times K^{\bullet}(X)
\xrightarrow{v}K^{\bullet}(X')\longrightarrow0.
\]

Tout d'abord, \(u\) admet l'inverse à gauche \(u'\), et a fortiori est injectif. En effet, on a:
\[
g_*\bigl(g^*(y).\lambda_{-1}(F)\bigr)
=y.g_*(\lambda_{-1}(F))
\]
d'après la formule de projection. Par ailleurs, \(g_*(\lambda_{-1}(F))=1\) (VI 5.9) d'où le résultat.

Ensuite, \(v\circ u=0\). Si \(y\in K^{\bullet}(Y)\), on a en effet:
\[
v\circ u(y)
=v\bigl(g^*(y)\lambda_{-1}(F),-i_*(y)\bigr)
=j_*\bigl(g^*(y)\lambda_{-1}(F)\bigr)-f^*(i_*(x))=0,
\]
d'après 3.4.

Il reste donc à montrer que \(\Ker(v)\subset \Im(u)\). Soient donc \(x\in K^{\bullet}(X)\), \(y'\in K^{\bullet}(Y')\), tels que:
